\section{Why the numerical operations are efficient}
\label{sec:local}

Weakly order-preserving maps are classical counting objects
\citep{Stanley1986}. Here the intervals and recursive admissions let us use
them without counting the full public partial order.
We need comparison probabilities, cumulative coordinate probabilities, and
a reusable interval of positions with adequate point probabilities.

\begin{lemma}[Central mass and anti-concentration]
\label{lem:insertion}
Suppose $x$ has effective interval $[\ell,r]$, width $W$, and at most $c$
other active predecessors and successors in total. Let
$t_0=\lfloor(W-1)/4\rfloor$ and $I_x=[\ell+t_0,r-t_0]$. Then
$|I_x|\geq W/2$ and
\begin{align}
\Prob(X_x=t)&\geq \frac{1}{4^cW} &&(t\in I_x),\label{eq:central}\\
\Prob(X_x\in I)&\leq\frac{2\cdot4^c|I|}{W}
 &&\text{for every integer interval }I.\label{eq:concentration}
\end{align}
For a different coordinate $Y$, even a dependent one, or a fixed real
threshold $Y$,
\begin{equation}
\Prob(|X_x-Y|\leq z)\leq \frac{2\cdot4^c(2z+1)}{W}
\qquad(z\geq0).
\label{eq:near-tie}
\end{equation}
\end{lemma}

\begin{lemma}[Local approximation]
\label{lem:local}
Fix degree and height bounds $d,h$, a component width ratio $R\geq1$,
and an error $0<\varepsilon<1/4$. Explicit computable constants
$W_0,\rho$ have the following property. If every nontrivial active
component has minimum width at least $W_0$ and maximum-to-minimum width
ratio at most $R$, restricting to radius $\rho$ around one or two requested
coordinates, with their original propagated intervals, changes each requested
weak comparison or cumulative probability by at most $\varepsilon$.
Isolated coordinates can be treated exactly at any width.
\end{lemma}

The proof in Appendix~\ref{app:local} couples two conditional-refresh
experiments, one on the full graph and one on the truncated graph.
These experiments are used only in the proof. The adversary's computation
is deterministic and uses the following exact finite counter.

\subsection{The subset counter}
\label{sec:counter}

Let the restricted graph have vertex set $V$, $b=|V|$, and positions
$1,\ldots,K$. A set is downward closed if it contains every predecessor
of each of its members. In an admissible assignment,
$S_t=\{x:X_x\leq t\}$ is downward closed, and $S_{t-1}\subseteq S_t$.

Initialize $C_0(\varnothing)=1$ and the other entries to zero.
For each $t=1,\ldots,K$, and downward-closed $S\subseteq U\subseteq V$,
perform
\begin{equation}
                         C_t(U)\mathrel{+}=C_{t-1}(S)
\label{eq:dp}
\end{equation}
if the interval of every $x\in U\setminus S$ contains $t$.
This assigns exactly those new items position $t$. No internal ordering
is needed, since equal positions are permitted. The chain of sets $S_t$
recovers the assignment uniquely; thus $C_K(V)$ is exactly the count.

There are at most $3^b$ pairs $S\subseteq U$: each item is outside $U$,
in $U\setminus S$, or in $S$. Precomputing closure and transition conditions
takes a polynomial factor in $b$. The counter uses at most $K3^b$
transitions on integers of $O(b\log K)$ bits. A comparison numerator adds
the requested weak edge; a cumulative numerator adds an upper bound.
Even if the added edge creates a weak directed cycle, downward closure
forces all its vertices to enter together, correctly counting equality.

Use a single neighborhood for all cumulative thresholds. The resulting
estimates then form the cumulative distribution of an actual local law and
are automatically monotone. For a fixed endpoint $c$, the two answer
probabilities are $F_x(c)$ and $1-F_x(c-1)$.
Since the two true weak probabilities sum to at least one, choosing the
larger estimate retains at least $1/2-2\eta$. We use this conservative
bound throughout.

\begin{lemma}[An interval from approximate quartiles]
\label{lem:quartiles}
Let $\widehat F$ be a local coordinate cumulative distribution with pointwise
error at most $1/16$. Let $a,b$ be its first crossings of $1/4,3/4$.
Then $[a,b]$ lies in the true effective support and
\[
                   \Prob(X_x=t)\geq\frac{1}{16(b-a+1)}
                   \qquad(a\leq t\leq b).
\]
\end{lemma}

Appendix~\ref{app:quartiles} proves this using a counting injection establishing
log-concavity of coordinate probabilities. A likely interval is needed here,
rather than just a likely point: representative reuse is what controls the
total cost of moves.

\subsection{The width repairs meet the oracle hypotheses}

\begin{lemma}[Strict cuts and separated scales]
\label{lem:width-repairs}
Rule W2 retains at least $1-\eta$ of the old count, preserves two-position
supports, and implies the retired comparison after all future insertions.
In W3 the selected items stay free with unchanged effective intervals and
an incident edge until moved. Their representative choices retain a fraction
at least
\begin{equation}
                   \frac{1}{CVb(2+\logtwo(K/b))}
\label{eq:representative-cost}
\end{equation}
when chosen block size is $b$, with $V=2\cdot4^B$ and an absolute $C$
independent of $J$. At most two selected items enter any child.
\end{lemma}
\begin{proof}
If $W_u>LW_v$, endpoint monotonicity gives $r_u\leq r_v$.
The part of $u$'s interval rejected by $X_u\leq\ell_v-1$ contains at
most $W_v$ positions. By \eqref{eq:concentration} its probability is at most
$C_*W_v/W_u<\eta$. The retained interval has length at least
$W_u-W_v\geq2$. Changed predecessor upper bounds remain strictly above
their old lower bounds, since $\ell_v-1>\ell_u$.
The artificial intervals on opposite sides of the cut are disjoint and
ordered. Replacing a position by adjacent copies never destroys that strict
separation. The other width direction is symmetric.

When W2 is unavailable, adjacent widths have ratio at most $L$.
The item selected for scale $i$ has width in $[A^iw,LA^iw]$.
Consecutive selected widths have ratio at least $A/L>L^{2h}$.
Even with shortcuts in the graph, they are therefore at distance greater
than $2h$. Fixing an item affects only its active predecessors and
successors, all within distance $h$, so it changes neither another selected
item's interval nor its incident edges.

For each previous width at most $w'$, a previously selected block has size
at most $w'$. It forbids at most one block of the same kind in its batch.
Thus previous choices exclude at most $Jw'$ positions from the present
central interval. The next width is greater than $32Jw'$, while its central
interval contains at least half that width. After these exclusions and
endpoint trimming a constant fraction remains. The largest-block estimate
in Lemma~\ref{lem:dictionary} still applies with an absolute constant.
The two kinds of representative permit at most two admissions at a position.
\end{proof}

Width repairs remove an edge or a free item, and therefore terminate.
Afterwards every nontrivial component satisfies Lemma~\ref{lem:local}.
Adding one new edge to a graph of height less than $H$ joins at most two
old directed paths; its height is less than $2H$. The degree rises by at
most one. All ensuing repairs only remove vertices or edges. Thus the
bounds $d,h,B$ used above remain valid.
All expensive numerical queries occur after width repairs.
